\subsection{The full-mean exponential lemma and the minimum-two root catalogue}\label{proof:Q5-C081}
\paragraph{Q81.0 Local message bounds}\label{part:Q81:local-roles}
Use the finite skew-kernel model of Section~\ref{proof:Q5-C064} and put
\(W_n=A\bar f_{n-1}\). The root has distinct residual arm lengths
\(2,u_1,v_1,\ldots,u_j,v_j\), where
\[
2<u_1<v_1<\cdots<u_j<v_j,\qquad j\ge1.
\]
The terminal has residual arm lengths \((1,b)\), with \(b\ge2\).

For a unit connector, set \(P_{\rm unit}=\prod_{i=1}^jW_{u_i}\bar f_{v_i}\),
except that when \(b=2,j=1,u_1=5\), take
\(P_{\rm unit}=h_5f_{v_1}\). The message and bound are
\begin{equation}\tag{Q81.0.1}\label{eq:Q81:unit-role}
Y_{\rm unit}=f_2P_{\rm unit}T(p\bar f_b),\qquad
\E Y_{\rm unit}^8\le e^{-191d/24}.
\end{equation}

For a connector of length \(L\ge2\), restrict to \(u_1=3\) and write
\(v=v_1\). The message is
\begin{equation}\tag{Q81.0.2}\label{eq:Q81:nonunit-message}
Y_{\rm nonunit}=W_2\bar f_3W_v
\prod_{i=2}^jW_{u_i}\bar f_{v_i}\,T^L(qf_b).
\end{equation}
It satisfies \eqref{eq:Q81:nonunit-bound}. The first pair is
\(\bar f_3W_v\), and its direction is fixed by the arm lengths and L.

These bounds include arbitrary genuine equal-pair products at the terminal
and root. Insert the terminal product inside T or \(T^L\) and multiply
the root product outside. Each product lies in \([0,1]\), so positivity
bounds the decorated message by the displayed undecorated message.

\paragraph{Q81.1 Original finite calculus and finite masks}

Delete zero-mass atoms and use the measure on the remaining atoms. Since diagonal
kernel entries vanish, \(p(x),\allowbreak  q(x)\ge\mu(x)>0\). All finite tree
messages are positive: the all-equal labeling has positive mass and
diagonal edge factors one. Both A and T are positive operators.

Set

\fitdisplay{
r=Kg,\quad \sigma=Kr,\quad t=K\sigma,\quad
 \delta_i=\|K^{i}\one\|_2^2,\quad
 \theta=13/32,\quad \kappa=19/32,\quad \tau=d-\delta_2.
}{}

The C014 calculus supplies

\fitdisplay{
\|K\|^2\le\theta,\quad0\le d\le1/3,\quad
 d^2\le \delta_2\le\theta d,\quad \delta_{i+1}\le\theta \delta_i,
 \quad |KV(x)|\le\|V\|_2,
}{Q81.1.1}
\fitdisplay{
Er=-d,\quad E\sigma=0,\quad Et=\delta_2,\quad Egr=0,
 \qquad\kappa d\le\tau\le d-d^2\le2/9.
}{Q81.1.2}

The original two-row bound \(|K^2V(x)|^2\le\|V\|^2/3\) is used
only in quoting specified already-paid original scalar certificates.
It is not a spectral bound for an arbitrary compression.

Define the finite quantities

\fitdisplay{
G_n(z)=\sum_{i=0}^{n-1}z^i,\quad
 \beta_1=\beta_2=0,\quad
 \beta_n=\sum_{i=1}^{\lfloor(n-1)/2\rfloor}(-1)^{i-1}\delta_i.
}{}

For either sign the exact finite positive convolution is

\fitdisplay{
\sum_{i=0}^n(\pm K)^i1
 =f_n^{\pm}+\sum_{i=3}^n\beta_i f_{n-i}^{\pm},
 \qquad \tau\le\beta_n\le d\quad(n\ge3),
}{Q81.1.3}

where \(f_n^+=f_n,\allowbreak  f_n^-=\bar f_n\). This finite identity applies at every integer length.

For a rotation frequency write \(z=i\sqrt x\), \(0\le x\le\theta\),
and \(G_n(z)=a_n+z c_n\). For \(n\ge2\),

\fitdisplay{
1-x\le a_n,c_n\le1;\qquad a_1=1,\ c_1=0.
}{Q81.1.4}

For \(n\ge3\), \(a_n\le1-x+x^2\) and \(|G_n(z)|\le1\).
For \(n\ge4\), both coordinates are at most \(1-x+x^2\).
The finite geometric formula proves these statements: odd indices
have squared numerator \(1+x^n\), and even indices have
\((1\pm x^{n/2})^2\), all divided by \(1+x\). In the positive-ending
even case \(n\ge6\), use \(2x^2+x^5\le1\).
We will also use

\fitdisplay{
|G_n(z)|^2\le1-7x/15\qquad(n\ge3).
\label{eq:Q81:finite-geometric-loss}
}{Q81.1.5}

For \(n=3\), the left side is \(1-x+x^2\), and \(1-x\ge7/15\).
For \(n=4\), it is \((1-x^2)^2/(1+x)\le1/(1+x)\le1-7x/15\).
For every \(n\ge5\), the finite numerator is bounded by
\((1+x^3)^2\). The remaining comparison is equivalent to
\(8-7x-30x^2-15x^5\ge0\); this decreases on the interval and its
value at \(13/32\) is \(1311885/33554432>0\).
These arguments cover every integer tail.

The spectral measure of g is denoted by \(\nu\). It has mass d and
\(\int x^i\,\allowbreak  d\nu=\delta_{i+1}\). Apply the frequency inequalities through the spectral theorem for this kernel.

\paragraph{Q81.1.1 Actual extra pairs and their finite masks}

C067 gives \(h_u f_v\le e^{R_{uv}-3d}\) for every \(4\le u<v\),
where \(R_{uv}=\sum_{a=u}^{v-1}K^a g\). C014 gives the exact exceptional
loss \(2d-\delta_2\) at \(u=3\). Full reversal gives, for negative pairs,

\fitdisplay{
\prod_i W_{u_i}\bar f_{v_i}\le e^{S-\Gamma},\quad
 S=K U(-K)g,\quad
 U(z)=\sum_i\sum_{a=u_i}^{v_i-1}z^{a-1},
}{Q81.1.6}
\fitdisplay{
\Gamma=3jd\ (u_1\ge4),\qquad
 \Gamma=2d-\delta_2+3(j-1)d\ (u_1=3).
}{Q81.1.7}

The intervals are disjoint, so U is a finite zero-one mask. Separating
the four residue classes at \(i\sqrt x\) gives, for every finite
zero-one or [0, 1]-coefficient mask Q,

\fitdisplay{
|Q(i\sqrt x)|^2\le\frac{1+x}{(1-x^2)^2},\qquad
 \|Q(K)\|^2\le\frac{32768}{16245}<\frac{100}{49}.
}{Q81.1.8}

Geometric series bound the finite positive subcollections in these sums. Consequently

\fitdisplay{
\|U(-K)g\|\le
 \begin{cases}(3/8)\sqrt d,&u_1\ge4,\\
 (1/4)\sqrt d,&u_1\ge5,\\
 (3/5)\sqrt d,&u_1=3,
 \end{cases}
}{Q81.1.9}
\fitdisplay{
ES\le d/12\ (u_1\ge4),\qquad
 ES\le \delta_2+d/12\ (u_1=3).
}{Q81.1.10}

The positive mean terms after excluding the first quadratic term are
at most \(\delta_4/(1-\theta^2)<d/12\). For a single pair starting at three,
\(U=z^2G_h(z)\), \(h=v-3\ge1\), and \(ES\le \delta_2\).

When \(u_{1}=3\) the mandatory quadratic coefficient gives

\fitdisplay{
U=u-i\sqrt x\,w,\quad
 -6x/5\le u\le-x/2,\qquad -x/2\le w\le6x/5.
}{Q81.1.11}

For a single such pair, the stronger form is
\(U=-x(a_h+z c_h)\), where \(a_h\in[1-x,\allowbreak  1]\),
\(c_h\in[0,\allowbreak  1]\), and \(c_h\ge1-x\) if \(h\ge2\).

At \(d=0,\allowbreak  g=0\) on every retained atom. K need not be zero. Successive
leaf integration nevertheless makes every actual message in this
manuscript equal to one. All endpoint assertions follow directly.
Every subsequent division by d assumes \(d>0.\)

\paragraph{Q81.2 A signed-mean exponential lemma}\label{part:Q81.2}

Theorem T086.1. Let W be any finite real variable with
\(|W|\le U\) and \(EW^2\le\theta U^2\), where \(\theta=13/32\).
For every \(s\ge0\),

\fitdisplay{
\boxed{Ee^{sW}\le
 \exp\{sEW+(17/64)s^2U^2\}.}
}{Q81.2.1}

Consequently, for every real vector V in the original kernel model,

\fitdisplay{
\boxed{Ee^{8KV}\le e^{8E(KV)+17\|V\|_2^2}.}
}{Q81.2.2}

The case \(U=0\) is immediate. For \(U>0\), put \(m=EW\),
\(v=E(W-m)^2\), \(M=U-m\). The cases \(M=0\) or \(v=0\) follow
directly. Otherwise the quadratic Hermite interpolant tangent to
\(e^{sw}\) at \(m-v/M\) and agreeing at U majorizes it on \([-U,U]\):
its remainder has the sign of \((w-m+v/M)^2(w-U)\le0\).
The lower contact lies in the interval since \(v\le U^2-m^2\).
Matching the first two moments therefore gives

\begin{equation}\tag{Q81.2.3}\label{eq:Q102:3.3}
Ee^{sW}\le e^{sm}Ee^{sL(B-\pi)},\qquad
L=M+v/M,\qquad \pi=\frac{v}{M^2+v},
\end{equation}

where B is Bernoulli with mean pi. For \(m\ge0\),
\(L\le(1+\theta)U\). The centered two-point log moment has second
derivative at most \(L^2/4\); integrating twice gives the coefficient
\((1+\theta)^2/8<17/64\).

For \(m\le0\), put \(u=-m/U\in[0,\sqrt\theta]\). Increasing the raw
second moment to \(\theta U^2\) enlarges L and decreases \(y=1-2\pi\),
giving

\begin{equation}\tag{Q81.2.4}\label{eq:Q102:3.4}
\begin{gathered}
L/U\le\ell(u)=\frac{1+\theta+2u}{1+u},\qquad
y\ge y(u)=\frac{\kappa+2u+2u^2}{1+\theta+2u}\ge0.
\end{gathered}
\end{equation}

The Bernoulli bound is

\begin{equation}\tag{Q81.2.5}\label{eq:Q102:3.5}
Ee^{t(B-\pi)}\le
\exp\left\{\frac{1-2\pi}{4\log((1-\pi)/\pi)}t^2\right\}.
\end{equation}

Indeed, put \(a=\frac12\log((1-\pi)/\pi)\). The function
\(z\mapsto\log\cosh\sqrt z\) is concave on \(z\ge0\), since
\(\tanh a/a\) decreases; the derivative test is
\((\tanh a-a\operatorname{sech}^2a)'=2a\operatorname{sech}^2a\tanh a\ge0\).
Apply the tangent inequality at \(a^2\) to \((a-t/2)^2\) in
\(Ee^{t(B-\pi)}=e^{t\tanh(a)/2}\cosh(a-t/2)/\cosh a\).
The endpoints follow by continuity.

The positive-term integral for arctanh gives

\fitdisplay{
\log\frac{1+y}{1-y}\ge2(y+y^3/3),\qquad0\le y<1.
}{}

Thus the coefficient in (Q81.2.3) is at most
\(L^2/[8(1+y^2/3)]\). It remains to prove

\begin{equation}\tag{Q81.2.6}\label{eq:Q102:3.6}
8\ell(u)^2\le17(1+y(u)^2/3)
\qquad(0\le u\le\sqrt\theta).
\end{equation}

Set \(A_0=45+64u\) and \(B_0=19+64u+64u^2\). Then
\(\ell=A_0/[32(1+u)]\) and \(y(u)=B_0/A_0\). Clearing positive
denominators reduces (Q81.2.6) to \(\mathcal R(u)\ge0\), where

\fitdisplay{
\mathcal R(u)=2176(1+u)^2(3A_0^2+B_0^2)-3A_0^4.
}{}

Expansion gives the certificate

\begin{equation}\tag{Q81.2.7}\label{eq:Q102:3.7}
\begin{aligned}
\mathcal R(u)={}&1167597+918784u+1048576u^3+919552u^4\\
&+35651584u^5+8912896u^6
+8564224(2u^2-1/4)^2.
\end{aligned}
\end{equation}

Every term is nonnegative for \(u\ge0\), proving (Q81.2.1).
Applying the row and spectral bounds in (Q81.1.1) with
\(W=KV\), \(U=\|V\|_2\), and \(s=8\) gives (Q81.2.2).

\paragraph{Q81.3 The unit terminal and a complete correlated frequency estimate}

For \eqref{eq:Q81:unit-role}, use only the terminal comparisons of
Lemma~\ref{lem:ordinary-unit-inner-terminal}. For \(b\ge3\), take
\(\bar R,\allowbreak m_b,\allowbreak k_b,\allowbreak t_0,\allowbreak e,\allowbreak V_b,\allowbreak\gamma\) from
\eqref{eq:ordinary-unit-inner:terminal-data}, writing
\(m=m_b\), \(k=k_b\), and \(V=V_b\).
For \(b=2\), take \(m=k=1\), \(\gamma=d-\delta_2/2\), and
\(\sigma=K^2g\) in \eqref{eq:ordinary-unit-inner:terminal2}.
The terminal lemma has no restriction on the receiving length;
the receiving factor here remains the actual \(f_2\).

Since \(f_{2}=p+r,\allowbreak \) the arithmetic mean of the two positive factors yields

\fitdisplay{
f_2T(p\bar f_b)\le
 \begin{cases}
 [1-\gamma/2+K(1+z)g/2]^2,&b=2,\\
 m^2[1-\gamma/2+KQ(z)g/2]^2,&b\ge3,
 \end{cases}
 \quad z=-K,\quad Q(z)=1+kzG_{b-1}(z).
}{Q81.3.5}

The square bases are nonnegative actual arithmetic means. Thus their
sixteenth powers are bounded by the corresponding exponentials,
including any zero boundary by continuity.

\paragraph{Q81.3.1 Single pair starting at three, \(b\ge 3\)}

For every \(N\ge 2,\allowbreak  h\ge 1,\allowbreak \) k in [1, \(6/5\)] and x in [0, theta],

\fitdisplay{
\boxed{|1+kzG_N(z)+z^2G_h(z)|^2\le1-159x/160,
 \qquad z=i\sqrt x.}
}{Q81.3.6}

The real and normalized imaginary coordinates are
\(1-kx c_N-xa_h\) and \(ka_N-xc_h\). They are positive:
their lower bounds are \(17/160\) and \(3/16.\) It is therefore legitimate
to upper-bound each coordinate before squaring.

For \(h=1,\allowbreak  N=2,\allowbreak \) the squared norm minus one, divided by x, is

\fitdisplay{
k^2-2k-2+(k+1)^2x.
}{}

It is increasing in x and convex in k; at \(x=\theta\) its k endpoints
are -\(11/8\) and -\(159/160.\)
For \(h=1,\allowbreak  N\ge 3,\allowbreak \) the coordinate upper bounds are
\(1-(k+1)x+kx^2\), \(k(1-x+x^2)\). The same difference quotient is

\fitdisplay{
k^2-2k-2+(-k^2+4k+1)x+(k^2-2k)x^2-k^2x^3(1-x).
}{}

The final two terms are nonpositive. The remaining endpoint values
are -\(11/8\) and -\(951/800,\allowbreak \) both at most -\(159/160.\)
For \(h\ge 2,\allowbreak  N=2\) use coordinates
\(1-(k+1)x+x^2\), \(k-x+x^2\). The quotient is

\fitdisplay{
k^2-2k-2+(k^2+3)x-x^2-x^3+x^4.
}{}

Its last three terms are nonpositive. The remaining endpoints are
-\(11/8\) and -\(37/32.\) Finally, for \(h\ge 2,\allowbreak  N\ge 3,\allowbreak \) use
\(1-(k+1)x+(k+1)x^2\), \(k-(k+1)x+(k+1)x^2\). The quotient is

\fitdisplay{
k^2-2k-2+(-k^2+2k+3)x+(k^2-1)x^2-(k+1)^2x^3(1-x).
}{}

Drop only its last nonpositive term. The remainder increases in x
and is convex in k; its endpoints are -\(11/8\) and -\(32733/25600,\allowbreak \)
both at most -\(159/160.\) This covers the four cases.

\paragraph{Q81.3.2 General masks for the other \(b\ge 3\) unit rows}\label{part:Q81.3.2}

If \(u_{1}=3\) and \(j\ge 2,\allowbreak \) (Q81.1.11) gives

\fitdisplay{
|Q+U|^2\le(1-35x/32)^2+x(6/5+x/2)^2<3/2.
}{Q81.3.7}

The coordinates remain positive. The excess over one is
\(-299x/400+(1225/1024+6/5)x^2+x^3/4\), at most
\(5/12+1/56<1/2\), using \(x^2<1/6,\allowbreak  x^3<1/14.\)

For every mask starting in degree at least three,

\fitdisplay{
\boxed{|Q+U|^2\le4/3.}
}{Q81.3.8}

To prove it, put \(\delta=x^2/(1-x^2)\le 6x^2/5.\) The positive upper tail
in each coordinate is at most delta. Both total coordinates remain
positive; the possible negative tails have magnitudes at most
\(x^3/(1-x^2)\) and \(x/(1-x^2),\allowbreak \) respectively. For \(N=2\) the upper coordinates
are \(1-kx+6x^2/5\), \(k+6x^2/5\). Their squared norm is

\fitdisplay{
1+(k^2-2k)x+(k^2+12/5)x^2+(36/25)(x^4+x^5).
}{}

For \(N\ge 3\) they are \(1-kx+(k+6/5)x^2\) and
\(k-kx+(k+6/5)x^2\). Their squared norm is

\fitdisplay{
1+(k^2-2k)x+(2k+12/5-k^2)x^2+k^2x^3
 +(k+6/5)(6/5-k)x^4+(k+6/5)^2x^5.
}{}

Both are convex in \(x\ge 0\) and in k on the specified interval.
Thus check \(x=0,\allowbreak \) theta and \(k=1,\allowbreak  6/5.\) At \(x=\theta,\allowbreak \) with denominator
838860800, the four numerators are respectively
1015054740, 1089601940, 1080009045, 1111336016.
Each is below four thirds of that denominator; for the largest,
the cross-multiplication difference is \(21435152>0.\)
Convexity extends these endpoint bounds to the interval.

\paragraph{Q81.4 The unit-connector estimate}\label{part:Q81:unit-estimate}

For \(b\ge 3\) combine (Q81.3.5) and (Q81.1.6) before using (Q81.2.2). In every case,

\fitdisplay{
\log EY^8\le
 -8(\beta_b+\gamma-ke+d+\Gamma-ES)
       +17\|(Q+U)g\|^2.
}{Q81.4.1}

For \(u_1=3\), \(j=1\), the last line of \eqref{eq:ordinary-unit-inner:terminal-long} and (Q81.3.6) make the right side at most

\fitdisplay{
-23d+(23093/800)\delta_2\le-10d.
}{}

For \(u_{1}=3,\allowbreak  j\ge 2\) use \(\Gamma\ge 5d-\delta_2,\allowbreak  ES\le \delta_2+d/12,\allowbreak \) and (Q81.3.7).
The upper divided by d is at most
\(-64+(1144/25)\theta+2/3+51/2<-10\).
For \(u_{1}\ge 4\) use \(\Gamma\ge 3d,\allowbreak  ES\le d/12,\allowbreak \) and (Q81.3.8). The upper is
\((-48+(744/25)\theta+2/3+68/3)d<-10d\).
Further pairs only strengthen these losses. Thus every \(b\ge 3\) row
has \(EY^8\le e^{-10d}\).

For \(b=2\) the common logarithmic expression is

\fitdisplay{
-8\gamma-8d-8\Gamma+8ES+17\|(1+z+U)g\|^2.
}{Q81.4.2}

For a single \(u_{1}=3\) pair, \(1+z+U=G_{h+2}\). By (Q81.1.5), (Q81.4.2) is at most
\(-15d+(181/15)\delta_2\le-(4847/480)d<-10d\).
For \(u_{1}=3,\allowbreak  j\ge 2,\allowbreak \) (Q81.1.11) gives

\fitdisplay{
|1+z+U|^2\le(1-x/2)^2+x(1+x/2)^2<5/4.
}{}

Use \(\Gamma\ge 5d-\delta_2\) and \(ES\le \delta_2+d/12\) in (Q81.4.2); the result is below -10d.
For \(u_{1}\ge 4,\allowbreak  j\ge 2,\allowbreak \) (Q81.1.9) and \(\sqrt{1+\theta}<6/5\) give
\(\|(1+z+U)g\|^2<(6/5+3/8)^2d<5d/2\).
\(\Gamma\ge 6d\) and \(ES\le d/12\) again make (Q81.4.2) below -10d.

It remains to pay every single pair with \(u\ge 4.\) At \(u=4\) keep its
negative direction, so \(U=z^3G_h(z)\). At \(u=5\) use the positive
exception in \hyperref[part:Q81:local-roles]{paragraph Q81.0}; its primitive is
\(U=z^4G_h(-z)\). In both cases the real coordinate of 1+z+U
lies in [0, \(1+x^2\)], its normalized imaginary coordinate in [0, 1],
and \(ES\le 0.\) Therefore

\fitdisplay{
|1+z+U|^2\le1+x+2x^2+x^4.
}{}

With \(\Gamma=3d,\allowbreak \) (Q81.4.2) is at most
\(-23d+21\delta_2+34\delta_3+17\delta_5\le-8d\).
The last comparison has surplus \(413535/1048576\) at \(x=\theta.\)

For every \(u\ge 6\) keep the negative direction. The single interval has
\(|U|\le x^{5/2}\sqrt{1+x}\), so

\fitdisplay{
|1+z+U|^2\le(1+x)(1+(2/3)x^2)^2<7/4.
}{Q81.4.3}

The upper is increasing; at theta the numerator and denominator are
130816125 and 75497472. Its comparison with \(7/4\) has positive
difference 1304451 in that denominator. Using \(ES\le d/12,\allowbreak \)
(Q81.4.2) is now at most

\fitdisplay{
\left[-40+119/4+2/3+4\theta\right]d=-191d/24.
}{}

The cases above prove \eqref{eq:Q81:unit-role} for \(b\ge 2,\allowbreak \) every first short arm,
every finite number of pairs, and every finite last arm.

\paragraph{Q81.5 The exchanged first pair for the nonunit family}

For \eqref{eq:Q81:nonunit-message}, exchange the first pair of arms; the preceding direction has factor \(W_3\bar f_v\).

Apply A, a positive operator, to the negative finite convolution for
\(\bar f_{v-1}\). Since

\fitdisplay{
A\sum_{i=0}^{v-1}(-K)^i1=2-\beta_v-\sum_{i=0}^v(-K)^i1,
}{}

retaining the real final term \(\beta_{v-1}W_1=\beta_{v-1}p\) gives

\fitdisplay{
W_v\le2-\beta_v-S_v^- -\beta_{v-1}p,
 \qquad S_v^-=\sum_{i=0}^v(-K)^i1.
}{Q81.5.1}

Also \(\bar f_3=S_3^- -d\). Set

\fitdisplay{
h=v-3\ge1,\quad \alpha=\beta_v,\quad\beta=\beta_{v-1},\quad
 U(z)=z^2G_h(z),\quad S=K U(-K)g,
 \quad\Gamma=d+\alpha+\beta.
}{Q81.5.2}

Here \(S_v^- -S_3^-=S\). Adding (Q81.5.1) to the b3 identity and using
AM--GM on two nonnegative actual factors proves

\fitdisplay{
\boxed{\bar f_3W_v\le\exp\{-S-\beta g-\Gamma\},\qquad
 \Gamma\ge3d-2\delta_2\ge35d/16.}
}{Q81.5.3}

Exchanging the first pair changes the sign of S and gives the displayed loss.

Write \(G_h=a_h+z c_h\), with the conventions of (Q81.1.4), and put
\(\zeta=\delta_2/d\) when \(d>0.\) Direct skew integration gives

\fitdisplay{
S=KU g,\quad ES=\int x a_h\,d\nu\in[\kappa \delta_2,\delta_2],
 \quad E(Ug)=-\int x c_h\,d\nu.
}{Q81.5.4}

For \(h=1,\allowbreak  \alpha=\beta=d\), \(ES=\delta_2\), \(E(Ug)=0\).
For \(h=2,\allowbreak  \beta=d\), \(ES=\delta_2\), \(E(Ug)=-\delta_2\).
For \(h\ge 3,\allowbreak \)

\fitdisplay{
\tau\le\beta\le d-\kappa \delta_2,\quad |G_h|\le1,
 \quad E(Ug)\le-\kappa \delta_2.
}{Q81.5.5}

The bound on beta is the finite alternating comparison
\(\beta_n\le d-\delta_2+\delta_3\le d-\kappa \delta_2\) for \(n\ge 5.\)
For all \(h\ge 2,\allowbreak  E(Ug)\le-\kappa \delta_2\), while
\(\tau\le\beta\le d\). We retain these same-variable correlations.

\paragraph{Q81.6 The base table and its full-mean conversion}

Put \(X=T^L(qf_b)\), \(Y_0=W_2X=(p-r)X\), \(L\ge 2.\)
The C071/C075/C076 actual positive masks have the form

\fitdisplay{
0<X\le C[Bq+Br+Z-a],\qquad
 Y_0\le CB\left(1+\frac{Z-a}{2B}\right)^2,
 \quad Z=KV.
}{Q81.6.1}

The base of the square is the arithmetic mean of actual nonnegative
upper factors. The table below records the finite envelopes and scalar bounds used in this proof.

For \(b=2\) set \(\gamma=d-\delta_2/2,\allowbreak  B_2=1-\gamma\).
For \(b\ge 3\) use the same actual positive terminal remainder

\fitdisplay{
R=KG_{b-1}(K)g,\quad m=1-\beta_b/2,\quad k=1/m,
 \quad e=d-\beta_{b+2},\quad \rho=E(rR),
}{}
\fitdisplay{
\gamma=k\beta_{b+1}-k^2ER^2/2,\quad a_1=ke,
 \quad B_*=\gamma+a_1.
}{Q81.6.2}

Here \(1\le k\le6/5\), \(0\le e\le \delta_2\),
\(\kappa \delta_2\le e,\allowbreak  \rho\le \delta_2\), \(\gamma\le kd\),
\(a_1\le39d/80\), and \(1-\gamma\ge17/25\).
At a frequency, writing \(G_{b-1}=a_N+z c_N\),
\(e=\int x c_N\,\allowbreak  d\nu\) and \(\rho=\int x a_N\,\allowbreak  d\nu\).
Define

\fitdisplay{
P_n(z)=G_n(z)+z^{n+1}\quad(b=2),\qquad
 P_n(z)=G_n(z)+kz^{n+1}G_{b-1}(-z)\quad(b\ge3).
}{Q81.6.3}


\begingroup\small
\setlength{\tabcolsep}{3pt}
\setlength{\LTleft}{0pt}\setlength{\LTright}{\fill}
\begin{longtable}{@{}>{\raggedright\arraybackslash}p{0.1200\linewidth}>{\raggedright\arraybackslash}p{0.2000\linewidth}>{\raggedright\arraybackslash}p{0.2000\linewidth}>{\raggedright\arraybackslash}p{0.2000\linewidth}>{\raggedright\arraybackslash}p{0.2000\linewidth}@{}}
\caption{Finite masks for \(X=T^L(qf_b)\), indexed by \((b,L)\).}\label{tab:parameters:q5-c081:1}\\
\toprule
\((b,L)\) & C & B & a & Actual Z \\
\midrule\endfirsthead
\toprule
\((b,L)\) & C & B & a & Actual Z \\
\midrule\endhead
\bottomrule\endfoot
2, 2 & 1; & B2; & 0; & \(K^2(\gamma 1+r)\). \\
\addlinespace[2pt]
2, 3 & 1; & B2; & \(\tau\); & \(-K^2(g+\sigma)\). \\
\addlinespace[2pt]
2, 4 & 1; & \(B2-\tau\); & d B2; & \(K^2[\tau 1-B2 g+r+t]\). \\
\addlinespace[2pt]
\(2,\allowbreak L\ge 5\) & 1; & \(1-2\tau\); & \(\tau\); & \(K[\tau+z P_{L-2}(z)]g\). \\
\addlinespace[2pt]
\(b\ge 3,\allowbreak 2\) & \(m^2\); & \(1-\gamma\); & a1; & \(K^2(\gamma 1+kR)\). \\
\addlinespace[2pt]
\(b\ge 3,\allowbreak 3\) & \(m^2\); & \(1-B*\); & \(d-k \rho\); & \(K[a_{1}+z P1(z)]g\). \\
\addlinespace[2pt]
\(b\ge 3,\allowbreak 4\) & \(m^2\); & \(1-41d/80\); & \(\kappa d\); & \(K[41d/80+z P2(z)]g\). \\
\addlinespace[2pt]
\(b\ge 3,\allowbreak L\ge 5\) & \(1-d/2\); & \(1-d/2\); & \(d/2\); & \(K[d/2+z P_{L-2}(z)]g\). \\
\end{longtable}
\endgroup


In this table \(z=-K\); displayed strings such as P1 designate the finite
polynomials in (Q81.6.3), not an unspecified path input.

The complete old scalar expressions are

\fitdisplay{
8\log(CB)-8a/B+8M/B+
 \begin{cases}6D_2/B^2,&Z=K^2V_2,\\
 (2025/128)D_1/B^2,&Z=KV_1,
 \end{cases}
}{Q81.6.4}

with \(M\ge\max(EZ,\allowbreak  0)\) and the original specified vector bounds \(D_i.\)
Every expression (Q81.6.4) is at most -d/10. For \(b\ge 3,\allowbreak  L=2\) it is at most
-d/8. The improved b2L2 calculation in Section~\ref{proof:Q5-C075}, \hyperref[part:Q75.3.2]{paragraph Q75.3.2}, explicitly
pays that original scalar, with \(D_2=\delta_2\) and
\(M=\delta_2-\gamma d\); it is not merely a moment theorem.
The finite positive-end-sign family \(b=2\) mod 4 is included by C071's
same-remainder d3 payment. Its sign is not discarded here.

For reference, the whole-domain normalizations for the four two-row
rows, in their order above, are

\fitdisplay{
\frac{D_2}{B^2d}\le
 \theta\frac{324}{169},\quad\frac{324}{169},\quad
 \frac{217}{81}<\frac{25}{9},\quad\frac{540}{289}<2.
}{Q81.6.5}

The one-row normalizations, in their order, are

\fitdisplay{
\frac{D_1}{B^2d}\le
 \frac{289}{256},\quad\frac{481}{400},\quad
 \frac{27}{32},\quad\frac{18}{25};
}{Q81.6.6}

all are at most \(5/4.\) The lower B bounds are respectively
\(13/18,\allowbreak  13/18,\allowbreak  1/2,\allowbreak  5/9,\allowbreak  17/25,\allowbreak  1/2,\allowbreak  4/5,\allowbreak  5/6.\)
For example the b2L4 normalization is
\((1+4/9)^2+16/27=217/81\), using \(B_2=B+\tau\).
For \(b\ge 3,\allowbreak \) L2, use \(\theta(1+\theta)<3/5\), \(k\le 6/5,\allowbreak \)
and \(B\ge 17/25.\) The b2 tail quotient
\((\theta-2\theta\kappa\tau+\tau^2)/(1-2\tau)^2\)
increases for \(0\le \tau\le 2/9\) and has endpoint \(289/256.\)
These are the full original normalizations from Q5-C076.

Set \(V_0=V/B\) in (Q81.6.1). Changing (Q81.6.4) to the exact full-mean
17 scalar (Q81.2.2) costs at most \((17\theta-6)D_2/B^2=(29/32)D_2/B^2\)
in a two-row case, and \((17-2025/128)D_1/B^2=(151/128)D_1/B^2\)
in a one-row case. Actual \(EZ\le M\) is retained. Therefore

\fitdisplay{
\mathcal E_0:=8\log(CB)-8a/B+8E(KV_0)+17\|V_0\|^2
 \le c_0d,\qquad \|V_0\|^2\le5d/4,
}{Q81.6.7}
\fitdisplay{
c_0=5/2\ (b=2),\qquad
 c_0=8/5\ (b\ge3,L=2),\qquad
 c_0=3/2\ (b\ge3,L\ge3).
}{Q81.6.8}

The arithmetic is
\(-1/10+(29/32)(25/9)=3481/1440<5/2\),
\(-1/8+(29/32)(540/289)=1813/1156<8/5\), and
\(-1/10+(151/128)(5/4)=3519/2560<3/2\).
The norm bound in (Q81.6.7) follows from (Q81.6.5) after multiplication by
theta and from (Q81.6.6). 

\paragraph{Q81.6.1 Finite coefficient facts used below}

For every finite \(n\ge 3\) in (Q81.6.3), writing \(P_n=A_P+zB_P\),

\fitdisplay{
1-x\le A_P,B_P\le1.
}{Q81.6.9}

This is the exact C075 finite coefficient lemma. One direct
proof separates n modulo four in the finite alternating sums
\(S_a=\sum_{i=0}^a(-x)^i\). For \(n=4a\) the two additions have signs
+,+ and magnitudes \(kx^{2a+1}c_N,\allowbreak  kx^{2a}a_N\); for \(n=4a+1\) they
have signs -,+ and magnitudes \(kx^{2a+1}a_N,\allowbreak  kx^{2a+1}c_N\);
for \(n=4a+2\) they have signs -,- and magnitudes
\(kx^{2a+2}c_N,\allowbreak  kx^{2a+1}a_N\); and for \(n=4a+3\) they have signs
+,- and magnitudes \(kx^{2a+2}a_N,\allowbreak  kx^{2a+2}c_N\).
Each fits inside the final alternating interval because \(kx<1\) and
\(kx^2<1.\) At \(n=3\) its second coordinate is directly at least 1-x.
The four residue classes cover every integer n. The b2 polynomial is included
by \(k=1,\allowbreak  N=1.\) For b2 also \(A_P\le1-x+x^2\).

The two short positive-terminal polynomials have exact coordinates

\fitdisplay{
P_1:\ A_P=1-kx a_N,\ B_P=kx c_N;
 \qquad P_2:\ A_P=1-kx^2c_N,\ B_P=1-kx a_N.
}{Q81.6.10}

In particular \(A_{P_1}\ge41/80\), \(A_{P_2}\ge4/5\).
In the later tail we may use the weaker bounds \(A_P\ge 1/2\) and
\(B_P\ge 3/8,\allowbreak \) as well as their upper bound one.

\paragraph{Q81.7 Single exchanged pair: terminal \(b=2,\allowbreak \) every \(L\ge 2\)}

Let \(Y_1=Y_0\bar f_3W_v\). Combine the two actual remainders before
applying (Q81.2.2). Define

\fitdisplay{
\Delta=\|V_0-Ug-\beta\one\|^2-\|V_0\|^2.
}{Q81.7.1}

Equations (Q81.5.3),(Q81.6.1) give the exact scalar certificate

\fitdisplay{
\log EY_1^8\le\mathcal E_0-8\Gamma+
                         17\Delta-8ES.
}{Q81.7.2}

We prove, throughout the full continuous original domain,

\fitdisplay{
\boxed{17\Delta-8ES\le7d\qquad(b=2,L\ge2).}
}{Q81.7.3}

Together with (Q81.6.8) and \(\Gamma\ge 35d/16\) this proves
\(EY_1^8\le e^{-8d}\).
In the computations below \(\zeta=\delta_2/d\in[d,\allowbreak  \theta]\).
Spectral Cauchy also gives \(\delta_3\ge \delta_2^2/d\).

\paragraph{Q81.7.1 \(L=2\)}

Here \(V_0=(\gamma+z^2)g/B\), \(B=B2.\) Put \(\lambda=\gamma/B.\)
Then \(0\le\lambda\le5/13\), because
\(\gamma\le d-d^2/2\le5/18\). Since \(EV0=0,\allowbreak \) the beta cross term is
\(\beta^2+2\beta E(Ug)\le\beta^2\). The frequency increment is

\fitdisplay{
x^2|G_h|^2-(2/B)x^2a_h+2\lambda x a_h.
}{Q81.7.4}

For \(h=1,\allowbreak \) this gives \(\Delta\le d^2+2\lambda \delta_2-\delta_3\).
For \(h=2,\allowbreak \) use \(|G_2|^2=1+x,\allowbreak  a_2=1\) to get
\(\Delta\le d^2+2\lambda \delta_2-\kappa \delta_3\).
As \(ES=\delta_2\) in both cases, their cost divided by d is at most

\fitdisplay{
17/3+(66/13)\zeta-17\kappa\zeta^2
 \le17/3+\frac{(66/13)^2}{68\kappa}<20/3<7.
}{Q81.7.5}

For the strict comparison use \(34848<54587.\) The h1 quadratic is
stronger. The concave quadratic bounds apply throughout the stated interval.

For \(h\ge 3,\allowbreak  |G_h|^2\le1\) and \(a_h\ge\kappa\) make the d3 term
in (Q81.7.4) nonpositive. By (Q81.5.5),

\fitdisplay{
\Delta\le2\lambda \delta_2+d^2(1-\kappa\zeta)^2.
}{}

Thus its cost divided by d is at most

\fitdisplay{
(170/13-8\kappa)\zeta+(17/3)(1-\kappa\zeta)^2.
}{}

This increases on [0, theta]: its derivative is at least
\(170/13-8\kappa-34\kappa/3>0\).
At theta it is \(20904801/3145728<7\). This proves the full L2 row.

\paragraph{Q81.7.2 \(L=3\)}

Here \(V_0=(z+z^3)g/B\), \(B=B2,\allowbreak \) and \(EV_0=\tau/B\).
Because \(\beta\le d\) and \(2\tau/B\ge 2\kappa d>d,\allowbreak \) the entire beta cross
\(\beta^2-2\beta EV_0+2\beta E(Ug)\) is nonpositive.
The remaining increment is at most

\fitdisplay{
\int x^2\{1+x+(2/B)(1-x)\}\,d\nu
 \le\int\frac{x^2(49-23x)}{13}\,d\nu.
}{Q81.7.6}

Use \(B\ge 13/18.\) After subtracting \(8ES\ge 8\kappa\) d2, the frequency
upper has its maximum at \(x=\theta\) or zero: its quotient by x is
increasing, since \(49-46x>0.\) At theta the value is
\(217217/32768<7\). This proves L3.

\paragraph{Q81.7.3 \(L=4\)}

Now \(B=B2-\tau\) and
\(V_0=[\tau+B_2z+z^2+z^4]g/B\), so \(EV_0=B_2d/B\).
The frequency increment excluding beta is

\fitdisplay{
x^2|G_h|^2-(2/B)x^2(1-x)a_h+2x^2c_h
                 +(2\tau/B)(xa_h+x^2c_h).
}{Q81.7.7}

The beta term is at most \(\tau^2-2\tau B_2d/B\): the function
\(\beta^2-2\beta B_2d/B\) decreases on [tau, d], and \(EUg\le 0.\)
Combining it with the last term of (Q81.7.7) gives a nonpositive quantity.
Indeed it suffices that

\fitdisplay{
2[B_2-(1+\theta)\zeta]-B(1-\zeta)\ge0.
}{}

Its exact left side is
\(1-(1+2\theta)\zeta-\zeta(5-3\zeta)d/2\).
It decreases in both variables on the original rectangle, and at
\(\zeta=\theta,\allowbreak  d=1/3\) equals \(47/6144>0.\)

In the remainder of (Q81.7.7), use \(B\le 1,\allowbreak  a_h\ge 1-x,\allowbreak  c_h\le 1,\allowbreak \) and \(|G_h|^2\le 1+x.\) This gives

\fitdisplay{
\Delta\le\int x^2(1+5x-2x^2)\,d\nu.
}{}

After subtraction of 8kappa d2 the frequency quotient by x increases,
and its endpoint value is \(2961647/524288<7\). This proves L4.

\paragraph{Q81.7.4 Every \(L\ge 5\)}

Here \(V_0=[\tau+zP_n]g/B\), \(n=L-2\ge 3,\allowbreak  B=1-2\tau\ge 5/9.\)
Write \(Pn=A+zC.\) In this b2 family,
\(1-x\le A,\allowbreak  C\le1\), \(A\le1-x+x^2\), and \(EV_0\ge\kappa d/B\).
For \(h=1,\allowbreak \) the frequency term is at most
\((1-2\kappa/B)\delta_3+(2\tau/B)\delta_2\).
With \(\beta=d,\allowbreak \) the constant term is at most \(d^2-2\kappa d^2/B\).
Since \(\tau \delta_2\le\kappa d^2\), these give \(\Delta\le d^2\).
Thus \(17\Delta-8ES\le9d^2\le3d\).

For \(h\ge 2,\allowbreak \) the frequency cross satisfies

\fitdisplay{
A c_h-C a_h\le(1-x+x^2)-(1-x)^2=x.
}{Q81.7.8}

The terms involving tau and beta together are at most

\fitdisplay{
\tau d\{1-2\kappa\zeta-(2/B)(\kappa-\zeta)\}
 \le(73/512)\tau d\le(73/2304)d.
}{Q81.7.9}

To see this, use \(\beta^2\le \beta\) d, \(\beta\ge \tau,\allowbreak  EUg\le -\kappa\) d2,
\(EV0\ge \kappa\) d/B, \(ES\le \delta_2,\allowbreak \) and \(B\le 1.\) The coefficient of beta
before this substitution is negative. Its final upper is
\(1-2\kappa+2\theta\zeta\le73/512\).

Equations (Q81.7.8)--(Q81.7.9) therefore give

\fitdisplay{
\Delta\le \delta_3+(23/5)\delta_4+(73/2304)d.
}{}

Consequently

\fitdisplay{
17\Delta-8ES\le
 \left[17\theta^2+(391/5)\theta^3+1241/2304-8\kappa\theta\right]d
 =\frac{9817163}{1474560}d<7d.
}{Q81.7.10}

Use \(d3\le \theta\) d2, \(d4\le \theta^2\delta_2,\allowbreak \) followed by \(\delta_2\le \theta\) d;
the coefficient being multiplied by d2 is positive. This proves
every integer \(L\ge 5\) and completes (Q81.7.3).

\paragraph{Q81.8 Single exchanged pair: every \(b\ge 3\) and every \(L\ge 2\)}

For these rows we prove

\fitdisplay{
\boxed{17\Delta-8ES\le8d.}
}{Q81.8.1}

The same exact definition (Q81.7.1) and the same actual terminal parameters
are used. Together with (Q81.6.8), this yields exp(-79d/10) at L2 and exp(-8d)
at every \(L\ge 3.\)

\paragraph{Q81.8.1 \(L=2\)}

Put \(c=k/B\) and \(\lambda=\gamma/B.\) Then
\(c\le30/17\), \(\lambda\le cd\), and

\fitdisplay{
V_0=[\gamma+kz^2G_N(-z)]g/B,\quad N=b-1\ge2,
 \qquad EV_0=c e\ge c\kappa \delta_2.
}{Q81.8.2}

Its frequency increment is

\fitdisplay{
x^2|G_h|^2+2\lambda x a_h
       -2c x^2(a_Na_h-xc_Nc_h).
}{Q81.8.3}

For \(h=1,\allowbreak  \beta=d\) and \(EUg=0.\) Combining the gamma and beta terms
gives at most \(d^2+2c\theta d \delta_2\), since \(e\ge \kappa\) d2.
The remaining frequency contribution is at most \(-3\delta_3/16\).
Thus

\fitdisplay{
17\Delta-8ES
 \le17d^2+60\theta d \delta_2-8\delta_2-(51/16)\delta_3
 \le(17/3)d+\delta_2/8<6d.
}{Q81.8.4}

For \(h=2,\allowbreak  \beta=d\) and \(EUg=-\delta_2.\) The gamma plus constant part is
at most \(d^2,\allowbreak \) because \(ctheta<1.\) The cross in (Q81.8.3) is nonnegative
before its minus sign: \(a_N-xc_N\ge1-2x>0\). Hence
\(\Delta\le d^2+\delta_3+\delta_4\), and

\fitdisplay{
17\Delta-8ES\le[17/3+\{17\theta(1+\theta)-8\}\theta]d<7d.
}{Q81.8.5}

For \(h\ge 3,\allowbreak  |Gh|\le 1\) and
\(a_Na_h-xc_Nc_h\ge1-3x+x^2\ge-55/1024\).
The part of (Q81.8.3) not involving gamma is at most
\((5177/4352)\delta_3<(6/5)\delta_3\).
The gamma and beta terms together are at most \(d^2.\) For this bound, Write \(\beta=t\) d,
\(1-\zeta\le t\le1-\kappa\zeta\). After using \(\lambda\le cd,\allowbreak  EUg\le -\kappa\) d2, and (Q81.8.2), their upper divided by \(d^2\) is

\fitdisplay{
t^2+(60/17)\zeta-(893/272)t\zeta.
}{Q81.8.6}

It is convex in t. At \(t=1-\zeta\) its excess over one is
\(\zeta[-477+1165\zeta]/272\le0\).
At \(t=1-\kappa\) zeta that excess is
\(\zeta[-16/17+(40071/17408)\zeta]\le0\).
At theta the two brackets are strictly negative, with numerators
-\(119/32\) and -\(3365/557056,\allowbreak \) respectively. Thus (\(Q81.8.6)\le 1\) throughout.
It follows that

\fitdisplay{
17\Delta-8ES\le
 [17/3+(102/5)\theta^2-8\kappa\theta]d
 =\frac{54557}{7680}d<8d.
}{Q81.8.7}

This proves every \(b\ge 3\) terminal, including all finite end signs, at L2.

\paragraph{Q81.8.2 \(L=3\)}

Here \(V_0=[\delta+zP_1]g/B\), \(\delta=ke,\allowbreak  B\ge 1/2.\)
Use (Q81.6.10). The general frequency increment, before beta, is

\fitdisplay{
x^2|G_h|^2+(2\delta/B)x a_h
            +(2/B)x^2(A_Pc_h-B_Pa_h).
}{Q81.8.8}

For \(h\ge 2,\allowbreak \) all terms involving delta and beta have nonpositive sum.
To verify this fully, write \(\beta=t\) d, \(1-\zeta\le t\le 1.\) Since
\(EV_0=(d-k\rho)/B\ge(d-k\delta_2)/B\) and \(\delta\le k\) d2, their upper
divided by \(d^2\) is

\fitdisplay{
t^2-2\kappa t\zeta-(2/B)[t(1-k\zeta)-k\zeta^2].
}{}

The bracket is at least \(1-(1+k)\zeta\ge17/160>0\).
Since \(B\le 1,\allowbreak \) the last display is at most

\fitdisplay{
t^2-2t+2(k-\kappa)t\zeta+2k\zeta^2.
}{Q81.8.9}

At \(t=1\) it is \(-1+2(k-\kappa)\zeta+2k\zeta^2\);
at \(t=1-\zeta\) it is \(-1+2(k-\kappa)\zeta+(1+2\kappa)\zeta^2\).
Both increase in k, zeta and are negative at \(k=6/5,\allowbreak  \zeta=\theta.\)
Convexity in t proves the assertion on the full interval.

For \(h=1\) the same terms are at most \(d^2,\allowbreak \) since
\(1-k\zeta(1+\zeta)>0\); the remaining cross in (Q81.8.8) is
nonpositive. Thus \(\Delta\le d^2+\delta_3\), which pays (Q81.8.1).

For \(h=2,\allowbreak  A_P-B_P\le1-2\kappa x\). Using \(2/B\le 4\) gives
\(\Delta\le5\delta_3-15\delta_4/4\). As \(ES=\delta_2,\allowbreak \)

\fitdisplay{
17\Delta-8ES\le(852501/131072)d<8d.
}{Q81.8.10}

For \(h\ge 3,\allowbreak  |G_h|^2\le1\), and
\(A_Pc_h-B_Pa_h\le1-\kappa(1+\kappa)x\).
Hence \(\Delta\le5\delta_3-(969/256)\delta_4\), and

\fitdisplay{
17\Delta-8ES\le(65299507/8388608)d<8d.
}{Q81.8.11}

For both endpoint comparisons, the frequency upper divided by x is
increasing on [0, theta]. Its derivative is positive there:
respectively \(85-(255/2)x>0\) and
\(85-(16473/128)x>0\). The estimate follows for every frequency and h.

\paragraph{Q81.8.3 \(L=4\)}

Now \(V_0=[\delta+zP_2]g/B\), \(\delta=41d/80,\allowbreak  B=1-41d/80\ge199/240>4/5\).
The delta and beta terms in (Q81.8.8) have nonpositive sum for every h.
For \(h=1,\allowbreak \) their upper divided by \(d^2\) is

\fitdisplay{
1-(2/B)[1-(k\theta+41/80)\zeta]\le-1+2\zeta<0.
}{}

For \(h\ge 2\) use \(t=\beta/d\) in [1-zeta, 1], \(EUg\le -\kappa\) d2, and
\(EV_0\ge(d-k\delta_3)/B\). The bracket
\(t(1-k\theta\zeta)-(41/80)\zeta\) is positive: at its smallest
permitted t it is at least \(1-2\zeta+(39/80)\zeta^2>0\).
Replace 1/B by one only in its negative contribution. The resulting
upper is

\fitdisplay{
t^2-2t+2(k\theta-\kappa)t\zeta+(41/40)\zeta
 \le-\kappa+(41/40)\theta<0.
}{Q81.8.12}

For \(h=1\) the remaining frequency cross is nonpositive, so (Q81.8.1)
is immediate. For \(h\ge 2,\allowbreak \) the exact coordinates in (Q81.6.10) give

\fitdisplay{
A_Pc_h-B_Pa_h
 \le x[11/5-(153/80)x].
}{Q81.8.13}

Indeed its upper is \(A_P-B_P(1-x)\), then use \(a_N\le 1,\allowbreak  c_N\ge \kappa,\allowbreak \) and maximize the resulting expression in k.
The right side is nonnegative. With \(2/B\le 5/2\) and \(|Gh|^2\le 1+x,\allowbreak \)

\fitdisplay{
\Delta\le \delta_3+(13/2)\delta_4-(153/32)\delta_5.
}{}

Subtracting \(8ES\ge 8\kappa\) d2 gives an increasing frequency quotient
by x and endpoint value

\fitdisplay{
17[\theta^2+(13/2)\theta^3-(153/32)\theta^4]
       -8\kappa\theta
 =203700679/33554432<8.
}{Q81.8.14}

The derivative test is \(1+13x-(459/32)x^2>0\). This proves L4.

\paragraph{Q81.8.4 Every \(L\ge 5\)}

Here \(B=1-d/2\ge 5/6,\allowbreak  \delta=d/2,\allowbreak \) and \(V_0=[\delta+zP_n]g/B\),
\(n=L-2\ge 3.\) Use \(A_P\ge 1/2,\allowbreak  B_P\ge 3/8,\allowbreak \) and \(A_P,\allowbreak  B_P\le 1.\)
For \(h\ge 2\) the delta and beta terms have nonpositive sum. Their upper
divided by \(d^2\) is

\fitdisplay{
t^2-2\kappa t\zeta-(t-\zeta)/B
 \le t^2-t+\zeta-2\kappa t\zeta.
}{Q81.8.15}

Here \(t\ge 1-\zeta>\zeta.\) At \(t=1\) its value is (\(1-2\kappa)\zeta\le 0\);
at \(t=1-\zeta\) it is \(\zeta[(1+2\kappa)\zeta-2\kappa]\le0\).
Convexity gives the bound between these endpoints.

For \(h=1\) that sum is at most \((6/5)\zeta d^2\), and the frequency
part is at most d3. Therefore its cost is at most
\((913/160)\zeta d<8d\).

For \(h\ge 2,\allowbreak  A_Pc_h-B_Pa_h\le5/8+3x/8\). Since \(2/B\le 12/5,\allowbreak \)

\fitdisplay{
\Delta\le(5/2)\delta_3+(19/10)\delta_4,
}{}
\fitdisplay{
17\Delta-8ES\le
 [ (85/2)\theta^2+(323/10)\theta^3-8\kappa\theta]d
 =2375711d/327680<8d.
}{Q81.8.16}

The finite coefficient lemma gives these bounds for every \(n\ge 3\). This completes
(Q81.8.1) and every terminal \(b\ge 3,\allowbreak \) inner \(L\ge 2\) single-pair case.

\paragraph{Q81.9 Further pairs in the nonunit message}

The single-pair proof gives a scalar bound for the combined remainder. Its primitive is
\(V_*=V_0-Ug-\beta\one\). From (Q81.6.7), (Q81.1.4), and \(\beta\le d,\allowbreak \)

\fitdisplay{
\|V_*\|\le(9/8+1/2+3/5)\sqrt d
             =(89/40)\sqrt d<(9/4)\sqrt d.
}{Q81.9.1}

For the middle term use
\(\|Ug\|\le\theta\sqrt{1+\theta}\sqrt d<\sqrt d/2\);
for the last use \(\sqrt d<3/5\). These are original norm bounds.

All remaining pairs satisfy \(u_2\ge v+1\ge5\). Their complete
disjoint mask has primitive Vtail of norm at most \(\sqrt{d}\)/4, mean
\(E(KV_{tail})\le d/12\), and actual total loss \(3(j-1)d\).
Combining it with V* before (Q81.2.2) adds at most

\fitdisplay{
\left[\frac23+17\left(2\cdot\frac94\cdot\frac14+\frac1{16}\right)
                  -24(j-1)\right]d
 =\left[\frac{1001}{48}-24(j-1)\right]d.
}{Q81.9.2}

For \(j\ge 2\) this is at most -151d/48. The bound is uniform in the number, gaps, and largest length of the remaining arms. Combining (Q81.7.2)--(Q81.8.16) with (Q81.9.2) yields

\fitdisplay{
\boxed{E\left[W_2\bar f_3W_v\prod_{i=2}^jW_{u_i}\bar f_{v_i}
                         T^L(qf_b)\right]^8\le e^{-79d/10}
 \quad(b\ge2,L\ge2,j\ge1).}\label{eq:Q81:nonunit-bound}
}{Q81.9.3}

This proves \eqref{eq:Q81:nonunit-bound}, including every finite short connector and
both infinite integer tails. All continuous payments used only the
original domain \(d\le 1/3,\allowbreak  d\le \zeta\le 13/32,\allowbreak \) and the stated same-kernel
spectral moments. 
